\documentclass[11pt]{article}
\def\activityid{F09-F-v3}\usepackage[letterpaper,margin=.65in,top=.60in,bottom=.65in]{geometry}
\usepackage[T1]{fontenc}
\usepackage[default]{sourcesanspro}
\usepackage{microtype,tikz,array,tabularx,fancyhdr,amsmath,amssymb}
\usepackage[hidelinks]{hyperref}
\usetikzlibrary{calc,arrows.meta}
\setlength{\parindent}{0pt}\setlength{\parskip}{7pt}
\setlength{\headheight}{0pt}\setlength{\footskip}{26pt}\setlength{\emergencystretch}{2em}
\pagestyle{fancy}\fancyhf{}\renewcommand{\headrulewidth}{0pt}\renewcommand{\footrulewidth}{.4pt}
\fancyfoot[L]{\scriptsize Bellingham Math Circle / Week 9 / \activityid}
\fancyfoot[R]{\scriptsize \thepage}
\definecolor{ink}{gray}{.12}\definecolor{soft}{gray}{.95}\color{ink}
\newcommand{\PageTitle}[2]{{\fontsize{10}{12}\selectfont\bfseries WEEK 9\quad BOUNCE LAB\hfill #1\par}\vspace{3pt}{\fontsize{26}{29}\selectfont\bfseries #2\par}\vspace{2pt}}
\newcommand{\Subhead}[1]{\par\vspace{3pt}{\large\bfseries #1}\par}
\newcommand{\RuleBox}[1]{\par\begingroup\setlength{\fboxsep}{8pt}\colorbox{soft}{\parbox{\dimexpr\linewidth-16pt\relax}{#1}}\endgroup\par}
\newcommand{\WriteBox}[2]{\par\textbf{#1}\par\vspace{2pt}\begin{tikzpicture}\draw[black!40,line width=.5pt] (0,0) rectangle (\dimexpr\linewidth-.5pt\relax,#2);\end{tikzpicture}\par}
\newcommand{\RookBoard}[3]{\begin{tikzpicture}[x=#3,y=#3]\draw[step=1,black!45] (0,0) grid (#1,#2);\draw[thick] (0,0) rectangle (#1,#2);\node[font=\Large] at ({#1-.5},.5){$\star$};\draw[-{Latex},thick] (0,-.35)--(#1,-.35);\draw[-{Latex},thick] ({#1+.35},#2)--({#1+.35},0);\end{tikzpicture}}
\newcommand{\Table}[3]{\begin{tikzpicture}[x=#3,y=#3]\draw[step=1,black!25] (0,0) grid (#1,#2);\draw[thick] (0,0) rectangle (#1,#2);\fill (0,0)circle(3pt);\draw[-{Latex},very thick] (0,0)--(.7,.7);\node[below] at ({#1/2},-.1){width #1};\node[rotate=90,above] at (-.2,{#2/2}){height #2};\end{tikzpicture}}
\newcommand{\GraphNodes}{\coordinate(A)at(0,0);\coordinate(B)at(3,0);\coordinate(C)at(1.5,2.6);\coordinate(D)at(1.5,.95);}
\newcommand{\Kfour}[1]{\begin{tikzpicture}[scale=#1]\GraphNodes\draw[very thick](A)--(B)--(C)--(A)--(D)--(B) (D)--(C);\foreach \n in {A,B,C,D}{\node[circle,draw,fill=white,minimum size=22pt,font=\bfseries]at(\n){\n};}\end{tikzpicture}}
\newcommand{\House}[1]{\begin{tikzpicture}[scale=#1]\coordinate(A)at(0,0);\coordinate(B)at(3,0);\coordinate(C)at(3,2);\coordinate(D)at(0,2);\coordinate(E)at(1.5,3.3);\draw[very thick](A)--(B)--(C)--(D)--(A) (D)--(E)--(C);\foreach \n in {A,B,C,D,E}{\node[circle,draw,fill=white,minimum size=22pt,font=\bfseries]at(\n){\n};}\end{tikzpicture}}




\newcommand{\StudentPage}[1]{{\fontsize{10}{12}\selectfont\bfseries Week 9 / Billiards lab / #1\par}\vspace{10pt}}
\newcommand{\Problem}[2]{\par\noindent\textbf{Problem #1:} #2\par}
\newcommand{\Space}[1]{\par\begin{tikzpicture}\draw[black!35] (0,0) rectangle (\dimexpr\linewidth-.5pt\relax,#1);\end{tikzpicture}\par}
\newcommand{\Piles}[1]{\begin{center}\begin{tikzpicture}[x=1in,y=1in]\foreach \x in {0,...,#1}{\draw[black!50,rounded corners] ({\x*2.05},0) rectangle ({\x*2.05+1.8},1.3);}\end{tikzpicture}\end{center}}

\newcommand{\Rooms}[5]{\begin{tikzpicture}[x=#5,y=#5]\draw[step=1,black!20](0,0)grid(#3,#4);\foreach\x in{0,#1,...,#3}{\draw[thick](\x,0)--(\x,#4);}\foreach\y in{0,#2,...,#4}{\draw[thick](0,\y)--(#3,\y);}\fill(0,0)circle(3pt);\end{tikzpicture}}
\newcommand{\PlainGrid}[3]{\begin{tikzpicture}[x=#3,y=#3]\draw[step=1,black!30](0,0)grid(#1,#2);\fill(0,0)circle(3pt);\end{tikzpicture}}

\begin{document}

\PageTitle{FACILITATOR / 1 OF 6}{Bounces before mirrored rooms}
\RuleBox{F09 v3, September 2026: prepared and unpiloted. Applying Week 1 feedback is a design inference, not evidence these billiards pages have been taught. Tracing comes before unfolding; unfolding gets its own concrete work before arithmetic formulas.}
\Subhead{Preparation and access}
Print page 1 at each child's starting level (3 K--1, 3 middle, 1 upper: 7 sheets), with continuation masters and the 3-page extra available. Copy K page 3 separately when folding; cut between its two panels first. Keep earlier traced boards for later comparisons. Bring counters, pencils, rulers, scrap grid paper, and scissors for an adult. About 10 minutes' preparation; no actual ball.

K: adult reading/tracing allowed; child chooses turns. Middle: count lengths and multiples through 12, then count whole rooms. Upper: multiples, division and parity; gcd/lcm and the coprime lemma are optional adult-supported continuations. Extra: fractions and coordinates; supply irrationality if necessary. Oral reasoning and adult scribing are valid.
\Subhead{Common launch before separate starts}
0--10: let children handle counters and trace grid diagonals freely. 10--15: gather at a visible $2\times4$ table. Move diagonally from the black dot to the right wall. Ask a child to show the next move that stays inside and still heads upward. Dot that bounce and circle the ending corner. Show that interior line crossings do nothing. Everyone now knows the physical action and what to record; distribute different starts only afterward. Save mirror rooms for a later question about predicting a path.

15--35: trace and compare. 35--40: movement/reset. 40--55: revisit a puzzling path, fold one copy, or continue the room-count investigation. 55--60: share one checked prediction and tidy. These times are adjustable. Parent mainly supports K; organizer alternates middle/upper visits and leaves a specific next rectangle. Rotate mover, tracer, checker, so each child acts.
\Subhead{Hints and useful stopping points}
First check the bounce rule. Then enlarge or trace a difficult section while the child directs. Introduce mirrored rooms by physically folding a line across one wall, not by announcing the lcm. K may stop with contrasting corners. Middle may stop with matched bounce/crossing numbers, then later predict corners. Upper may stop with room-count parity and a verified bounce rule. General coprime division, no-bottom-left and complete inverse-design proofs remain adult conversations on pages 2--3. A returner can supply one familiar trace and proceed to the representation they have not used.
\newpage

\PageTitle{FACILITATOR / 2 OF 6}{Checked paths and corner rule}
\Subhead{K--1 and middle answers}
For width $w$, height $h$: $2\times2$ ends top-right with no bounce. $2\times4$ follows $(0,0),(2,2),(0,4)$, ending top-left; $4\times2$ follows $(0,0),(2,2),(4,0)$, ending bottom-right. Folding the two width-2, height-4 rooms turns the straight diagonal through $(2,2)$ into the first of these paths.

For $2\times3$, the folded path is $(0,0),(2,2),(1,3),(0,2),(2,0)$. In copied rooms, the first simultaneous wall meeting is $(6,6)$. The four middle rows have $(L,\text{widths},\text{heights},\text{corner})$:
\begin{center}\renewcommand{\arraystretch}{1.35}\begin{tabular}{c|c|c|c|c}Table&$L$&Widths&Heights&Corner\\\hline$2\times3$&6&3&2&bottom-right\\$3\times6$&6&2&1&top-left\\$3\times9$&9&3&1&top-right\\$4\times6$&12&3&2&bottom-right\end{tabular}\end{center}
A middle design reaching top-right after a bounce can be $2\times6$ (two bounces); top-left can be $2\times4$. Accept any verified design.
\Subhead{Upper theorem with all hypotheses}
Take positive integer $w,h$ and a 45-degree launch. Reflections unfold to $y=x$. A corner occurs when the common coordinate is both a multiple of $w$ and of $h$. Its first positive value is $L=\operatorname{lcm}(w,h)$. First use room counts $m=L/w,n=L/h$. If they shared $d>1$, distance $L/d$ would already cross $m/d$ widths and $n/d$ heights: an earlier corner. Thus they are coprime; both cannot be even, ruling out bottom-left without a gcd formula.

\textbf{Optional arithmetic bridge.} Write $w=ga,h=gb$ with coprime $a,b$. Any common multiple is $ga\,t$ with $b\mid at$. To justify $b\mid t$, begin with $a=2,b=3$: both $3t$ and $2t$ divide into threes, so their difference $t$ does too. In general, run repeated subtraction on strips of lengths $a,b$ and, alongside them, $at,bt$. Subtract the shorter from the longer until equal. Common divisors are unchanged, so the coprime base strips end at 1 and the scaled strips at $t$. Both scaled lengths remain multiples of $b$ throughout; hence $b\mid t$. The least positive $t$ is therefore $b$, and $L=gab$, $m=b,n=a$. The shared plan scaffold supplies physical bundle examples. If this step is supplied, record it as a given theorem.

An odd number of whole widths folds to the right, even to the left. An odd number of whole heights folds to the top, even to the bottom. Thus odd $a$, odd $b$ gives top-right; even $a$, odd $b$ gives bottom-right; odd $a$, even $b$ gives top-left. Both cannot be even, so first arrival at bottom-left is impossible. Scaling both dimensions preserves the reduced pair and corner.

Upper rows: $4\times6$: $L=12,m=3,n=2$, bottom-right; $6\times9$: $18,3,2$, bottom-right; $6\times10$: $30,5,3$, top-right; $8\times14$: $56,7,4$, bottom-right. Raw parity of the original lengths does not determine the endpoint.

\newpage

\PageTitle{FACILITATOR / 3 OF 6}{Bounces, inverse design, and slopes}
\Subhead{Count crossings, then design}
Before the final corner the line crosses $m-1$ vertical room walls and $n-1$ horizontal ones. With the optional lemma these are $b-1$ and $a-1$. No crossing belongs to both lists: that would be an earlier corner. So the number of bounces is $a+b-2$ (not counting launch or terminal corner). Path length, if wanted, is $\sqrt2 L$; $L$ itself is horizontal distance, not path length.

For top-right after exactly four bounces, $a,b$ are coprime odd positives with $a+b=6$. Candidates $(1,5),(3,3),(5,1)$ reduce to just $(1,5),(5,1)$ because $(3,3)$ is not coprime. Any positive integer scaling works, such as $2\times10$. Top-right after three bounces is impossible: two odd numbers have even sum, so $a+b-2$ is even. The latter is a parity obstruction, not a failed search.
\Subhead{Extra packet: rational and irrational directions}
The extra uses a \emph{unit square} and physical slope $s=p/q$ in lowest terms. A positive lattice point $(u,v)$ on $y=(p/q)x$ satisfies $qv=pu$, so its first occurrence is $(q,p)$. By the coprime lemma on page 2, $q$ divides $u$; the first positive choice is $u=q$, giving $v=p$. The endpoint's horizontal side is right iff $q$ is odd; its vertical side is top iff $p$ is odd. There are $p+q-2$ bounces. Slope $2/3$ ends bottom-right with 3 bounces; $3/2$ ends top-left with 3.

For slope $\sqrt2$, an integer corner $(u,v)$ with $u>0$ would give $\sqrt2=v/u$, impossible. If requested, prove irrationality: reduce a supposed fraction $v/u$ first. From $v^2=2u^2$, $v$ is even; substituting $v=2k$ forces $u$ even, contradicting the reduction.

The other launch $(0,1/2)$ with slope 1 cycles through $(1/2,1),(1,1/2),(1/2,0),(0,1/2)$. It avoids corners but follows just a diamond. No-corner does not imply dense. Density of the $\sqrt2$-slope trajectory launched from $(0,0)$ is a separate theorem. Under our stopping rule, arbitrary starts need care: from $(0,2-\sqrt2)$ the same slope reaches unfolded $(1,2)$, folds to bottom-right, and stops.
\Subhead{Keep the conventions straight}
A general rectangle with physical slope $s$ becomes a square after rescaling axes, with normalized slope $sw/h$. Rationality of that normalized slope controls the square model. The extra fixes $w=h=1$ to avoid silently mixing those slopes. ``Rational polygon'' in research means its angles are rational multiples of $\pi$; it does not mean a rational launch slope. Our corner-start paths stop when they hit a corner, so do not call those stopped paths periodic billiard orbits.

\newpage

\PageTitle{FACILITATOR / 4 OF 6}{A doorway into dynamics}
\Subhead{Mathematical direction and primary sources}
Unfolding replaces a piecewise reflected path by a straight line on a lattice. The corner problem is divisibility, coprime integers, and parity; the inverse problem asks for a complete classification of parameters. These are substantive undergraduate number-theoretic and geometric ideas in a tangible representation.

Masur and Tabachnikov, \emph{Rational billiards and flat structures}, published in \emph{Handbook of Dynamical Systems} 1A (2002), pp. 1015--1089. In the linked author manuscript, \S1.3 gives unfolding, \S1.4 develops the unit-square/torus model, and \S1.5 extends the construction to flat surfaces. Later \S4 concerns periodic orbits. \href{https://math.uchicago.edu/~masur/handbook1.pdf}{Author manuscript}.

The classroom genuinely performs the first construction used in the research: reflect copies, then study straight motion. It does not establish the deeper recurrence or ergodicity theorems. The extra explicitly separates avoiding corners from density; the latter is supplied specifically for the origin-start $\sqrt2$ path under our stopping convention.
\Subhead{Downloaded sources and documented prior use}
JRMF \emph{Billiards Geometry}, PDF p. 2, asks for corner, bounce-count, and path-length rules. Lowell Handout 10, problems 10.2--10.7, includes $7\times5$, $4\times3$, $8\times10$, scaling, consecutive sides, $1\times n$, and $2\times\text{odd}$ families. Handout 11, problems 11.4--11.5, includes height-five rectangles and plus-shaped tables. Do not describe small rectangle families or scaling as wholly new for returning children.

Our changed emphasis is a proof through physical unfolding, an exact first-corner criterion, a bounce formula, and inverse existence/impossibility design. This deliberately promotes an old reserve into the core. Plus-shaped tables remain a later branch; do not claim the rectangle formula covers them.

\emph{Math Circle by the Bay}, preface printed pp. viii--x, supports manipulatives and revisiting ideas at different depths. Rozhkovskaya, Lesson 3 ``At the lesson,'' EPUB \texttt{part0013.xhtml}, supports attempts and child explanations before formal displays; Lesson 8's account warns about copying pace. Preprinted grids and optional adult tracing are our response, not a procedure claimed to be in those sources.
\Subhead{Keep a useful history}
IDs F09-K/M/U/X-v3. Record actual dimensions, launch slopes, whether unfolding was used, which formulas were explained, and which claims stayed conjectural. A prepared extension is not a taught activity. Future variants can change slopes, starting points, or polygons instead of enlarging the same arithmetic table.

\texttt{verify.py} independently traces all 900 rectangles with sides 1--30, compares the corner and bounce formulas, and checks the inverse-design list. The general proofs are on pages 2--3; finite checks alone do not establish them.


\newpage
\PageTitle{FACILITATOR / 5 OF 6}{Representation gates and records}
\Subhead{Move, mark, compare; then unfold}
\textbf{K pages 1--2:} the square has no bounce, tall and wide tables turn at different wall types, and the height-6 table adds a second bounce. The child chooses directions even when an adult draws. Page 3 offers two separate folds, one across a vertical wall and one across a horizontal wall. Cut between the two diagrams; fold the right half left for Problem 5 and the top half down for Problem 6. A spare sheet or tracing paper makes the reflected segment visible. Folding is a follow-up to the paths, not an opening explanation.

\textbf{Middle:} page 1 gives three contrasting traces. Page 2 numbers bounces and unfolded crossings so the new picture has a concrete matching job. Demonstrate the first crossing (2,2); let the child match the rest. Page 3 gives separate 3-by-6 and 3-by-9 diagrams before translating parity into an endpoint table. Have the child actually move back and forth along a room edge as counts rise. Page 4 introduces lists of multiples to predict a larger table without drawing all copied rooms, then supplies design grids with stated bounds. No freehand rectangle is required to be metrically exact if the counted grid is correct.

\textbf{Upper:} page 1 contrasts a square and swapped side lengths. Page 2 gives original and copied 4-by-6 rooms side by side; the child matches actual bounces with crossings. Page 3 uses a scaled copy and two even-sided tables with different endpoints, preventing raw parity from being mistaken for the rule. The false distance 24 forces attention to the first corner. Page 4 counts separate wall lists before a formula. Page 5 compares three candidate designs, including the reducible 3-by-3 case. The two genuine four-bounce designs therefore contrast with an early-stop failure.
\Subhead{Preserve the proof as a conversation}
The least common distance argument, ``earlier corner if both room counts share a divisor,'' and wall-count proof can be shown directly on records. Only then offer reduced sides and the coprime-divisibility scaffold. If the child uses the formula but cannot justify its lemma, record it as supplied. The complete inverse-design proof and irrationality proof remain on pages 2--3. No child needs all these arguments in one hour.

The extra gives five rational trials, including 2/4 vs 1/2, before the general p/q claim. The midpoint diamond gives a concrete counterexample to ``no corner implies dense.'' Keep the origin hypothesis on the supplied irrational-slope density theorem. Arbitrary offsets and general rectangles need the qualifications on page 3.
\Subhead{After the meeting}
Record actual tables/slopes, bounce-rule misunderstandings, whether folding clarified the path, which records were useful, and questions children wanted to pursue. Distinguish an observed difficulty from a proposed explanation; record conjecture, explained proof, and supplied theorem separately. These v3 changes remain unpiloted until actual use is reported.
\newpage
\PageTitle{FACILITATOR / 6 OF 6}{Additional v3 solution checks}
\Subhead{K and middle}
K Problem 3: 4-by-2 ends bottom-right after 1 bounce; 2-by-6 follows (0,0),(2,2),(0,4),(2,6), ending top-right after 2. Problem 5 folds the straight 4-by-4 diagonal into the 2-by-4 path; Problem 6 folds it into the 4-by-2 path.

Middle page 1 corners/bounces: 2-by-2 top-right/0; 2-by-4 top-left/1; 4-by-2 bottom-right/1. Problem 4 unfolds 2-by-3 crossings at distances 2,3,4, ending at 6. The corresponding folded bounces are (2,2),(1,3),(0,2). Whole-room counts are 3 widths and 2 heights. Problem 6: 3-by-6 reaches (6,6), counts (2,1), top-left, 1 bounce; 3-by-9 reaches (9,9), counts (3,1), top-right, 2 bounces. Endpoint rows alternate right/left and top/bottom.

Problem 8: 4-by-6 has L=12, counts (3,2), bottom-right; 2-by-6 has L=6, counts (3,1), top-right; 3-by-3 has L=3, counts (1,1), top-right. Design examples within the 6-by-6 grids: 2-by-6 top-right after 2 bounces, 2-by-4 top-left after 1. A square fails the ``after at least one bounce'' condition.
\Subhead{Upper}
Problem 1: 2-by-3 ends bottom-right/3 bounces; 3-by-2 top-left/3; 3-by-3 top-right/0. Problem 4 for 4-by-6 has vertical distances 4,8,12 and horizontal 6,12: first common 12. The 6-by-9 scaling factor is 3/2; L scales 12 to 18 while counts and corner stay the same. 6-by-10 ends top-right; raw even side lengths do not determine corner. Problem 8's distance 24 predicts bottom-left, but the path stops at 12 first.

Problem 9: 2-by-3 vertical list 2,4, horizontal 3, total 3; square both empty, total 0; 4-by-6 lists 4,8 and 6, total 3; 6-by-10 lists 6,12,18,24 and 10,20, total 6. The count is (m-1)+(n-1). Doubling dimensions doubles all distances, preserving counts. Problem 12: 1-by-5 and 5-by-1 have four bounces and end top-right; 3-by-3 stops immediately. Problem 13 may use 2-by-10. For Problem 14, three bounces would give m+n=5; top-right requires both m,n odd, impossible. Page 3 gives the full reduced-pair classification.
\Subhead{Extra}
Slopes 1/2, 2/3, 3/2 first reach (2,1),(3,2),(2,3), with corners top-left, bottom-right, top-left and bounce counts 1,3,3. Slope 2/4 equals 1/2 and stops at (2,1), not (4,2). Slope 3/4 first reaches (4,3), top-left, 5 bounces. Problem 4 examples: 1/5 or 5 for top-right/4; 2/3 for bottom-right/3. For the midpoint diamond, (1/2,1/2) and a small neighborhood of it are never visited. Absence of corner hits alone does not establish density; the separate stated theorem does.

\end{document}
